% Options for packages loaded elsewhere
\PassOptionsToPackage{unicode}{hyperref}
\PassOptionsToPackage{hyphens}{url}
\documentclass[
]{article}
\usepackage{xcolor}
\usepackage[margin=1in]{geometry}
\usepackage{amsmath,amssymb}
\setcounter{secnumdepth}{-\maxdimen} % remove section numbering
\usepackage{iftex}
\ifPDFTeX
  \usepackage[T1]{fontenc}
  \usepackage[utf8]{inputenc}
  \usepackage{textcomp} % provide euro and other symbols
\else % if luatex or xetex
  \usepackage{unicode-math} % this also loads fontspec
  \defaultfontfeatures{Scale=MatchLowercase}
  \defaultfontfeatures[\rmfamily]{Ligatures=TeX,Scale=1}
\fi
\usepackage{lmodern}
\ifPDFTeX\else
  % xetex/luatex font selection
\fi
% Use upquote if available, for straight quotes in verbatim environments
\IfFileExists{upquote.sty}{\usepackage{upquote}}{}
\IfFileExists{microtype.sty}{% use microtype if available
  \usepackage[]{microtype}
  \UseMicrotypeSet[protrusion]{basicmath} % disable protrusion for tt fonts
}{}
\makeatletter
\@ifundefined{KOMAClassName}{% if non-KOMA class
  \IfFileExists{parskip.sty}{%
    \usepackage{parskip}
  }{% else
    \setlength{\parindent}{0pt}
    \setlength{\parskip}{6pt plus 2pt minus 1pt}}
}{% if KOMA class
  \KOMAoptions{parskip=half}}
\makeatother
\usepackage{color}
\usepackage{fancyvrb}
\newcommand{\VerbBar}{|}
\newcommand{\VERB}{\Verb[commandchars=\\\{\}]}
\DefineVerbatimEnvironment{Highlighting}{Verbatim}{commandchars=\\\{\}}
% Add ',fontsize=\small' for more characters per line
\usepackage{framed}
\definecolor{shadecolor}{RGB}{248,248,248}
\newenvironment{Shaded}{\begin{snugshade}}{\end{snugshade}}
\newcommand{\AlertTok}[1]{\textcolor[rgb]{0.94,0.16,0.16}{#1}}
\newcommand{\AnnotationTok}[1]{\textcolor[rgb]{0.56,0.35,0.01}{\textbf{\textit{#1}}}}
\newcommand{\AttributeTok}[1]{\textcolor[rgb]{0.13,0.29,0.53}{#1}}
\newcommand{\BaseNTok}[1]{\textcolor[rgb]{0.00,0.00,0.81}{#1}}
\newcommand{\BuiltInTok}[1]{#1}
\newcommand{\CharTok}[1]{\textcolor[rgb]{0.31,0.60,0.02}{#1}}
\newcommand{\CommentTok}[1]{\textcolor[rgb]{0.56,0.35,0.01}{\textit{#1}}}
\newcommand{\CommentVarTok}[1]{\textcolor[rgb]{0.56,0.35,0.01}{\textbf{\textit{#1}}}}
\newcommand{\ConstantTok}[1]{\textcolor[rgb]{0.56,0.35,0.01}{#1}}
\newcommand{\ControlFlowTok}[1]{\textcolor[rgb]{0.13,0.29,0.53}{\textbf{#1}}}
\newcommand{\DataTypeTok}[1]{\textcolor[rgb]{0.13,0.29,0.53}{#1}}
\newcommand{\DecValTok}[1]{\textcolor[rgb]{0.00,0.00,0.81}{#1}}
\newcommand{\DocumentationTok}[1]{\textcolor[rgb]{0.56,0.35,0.01}{\textbf{\textit{#1}}}}
\newcommand{\ErrorTok}[1]{\textcolor[rgb]{0.64,0.00,0.00}{\textbf{#1}}}
\newcommand{\ExtensionTok}[1]{#1}
\newcommand{\FloatTok}[1]{\textcolor[rgb]{0.00,0.00,0.81}{#1}}
\newcommand{\FunctionTok}[1]{\textcolor[rgb]{0.13,0.29,0.53}{\textbf{#1}}}
\newcommand{\ImportTok}[1]{#1}
\newcommand{\InformationTok}[1]{\textcolor[rgb]{0.56,0.35,0.01}{\textbf{\textit{#1}}}}
\newcommand{\KeywordTok}[1]{\textcolor[rgb]{0.13,0.29,0.53}{\textbf{#1}}}
\newcommand{\NormalTok}[1]{#1}
\newcommand{\OperatorTok}[1]{\textcolor[rgb]{0.81,0.36,0.00}{\textbf{#1}}}
\newcommand{\OtherTok}[1]{\textcolor[rgb]{0.56,0.35,0.01}{#1}}
\newcommand{\PreprocessorTok}[1]{\textcolor[rgb]{0.56,0.35,0.01}{\textit{#1}}}
\newcommand{\RegionMarkerTok}[1]{#1}
\newcommand{\SpecialCharTok}[1]{\textcolor[rgb]{0.81,0.36,0.00}{\textbf{#1}}}
\newcommand{\SpecialStringTok}[1]{\textcolor[rgb]{0.31,0.60,0.02}{#1}}
\newcommand{\StringTok}[1]{\textcolor[rgb]{0.31,0.60,0.02}{#1}}
\newcommand{\VariableTok}[1]{\textcolor[rgb]{0.00,0.00,0.00}{#1}}
\newcommand{\VerbatimStringTok}[1]{\textcolor[rgb]{0.31,0.60,0.02}{#1}}
\newcommand{\WarningTok}[1]{\textcolor[rgb]{0.56,0.35,0.01}{\textbf{\textit{#1}}}}
\usepackage{graphicx}
\makeatletter
\newsavebox\pandoc@box
\newcommand*\pandocbounded[1]{% scales image to fit in text height/width
  \sbox\pandoc@box{#1}%
  \Gscale@div\@tempa{\textheight}{\dimexpr\ht\pandoc@box+\dp\pandoc@box\relax}%
  \Gscale@div\@tempb{\linewidth}{\wd\pandoc@box}%
  \ifdim\@tempb\p@<\@tempa\p@\let\@tempa\@tempb\fi% select the smaller of both
  \ifdim\@tempa\p@<\p@\scalebox{\@tempa}{\usebox\pandoc@box}%
  \else\usebox{\pandoc@box}%
  \fi%
}
% Set default figure placement to htbp
\def\fps@figure{htbp}
\makeatother
\setlength{\emergencystretch}{3em} % prevent overfull lines
\providecommand{\tightlist}{%
  \setlength{\itemsep}{0pt}\setlength{\parskip}{0pt}}
\usepackage{bookmark}
\IfFileExists{xurl.sty}{\usepackage{xurl}}{} % add URL line breaks if available
\urlstyle{same}
\hypersetup{
  pdftitle={In-Class Exercise},
  pdfauthor={Yu-You Liou},
  hidelinks,
  pdfcreator={LaTeX via pandoc}}

\title{In-Class Exercise}
\author{Yu-You Liou}
\date{2026-06-02}

\begin{document}
\maketitle

\section{Role}\label{role}

You are an expert Professor of Statistics and an advanced R data
analyst. Your task is to programmatically and rigorously grade a
statistics assignment based on the provided directory structure, grading
rubrics, and reporting requirements.

\section{Variables \& Environment}\label{variables-environment}

\begin{Shaded}
\begin{Highlighting}[]
\SpecialCharTok{{-}} \StringTok{\textasciigrave{}}\AttributeTok{[no]}\StringTok{\textasciigrave{}} \OtherTok{=} \DecValTok{1}\NormalTok{ (representing Assignment }\DecValTok{1}\NormalTok{, e.g., ex1, ps1)}
\SpecialCharTok{{-}}\NormalTok{ Timezone}\SpecialCharTok{:}\NormalTok{ UTC}\SpecialCharTok{+}\DecValTok{8}
\SpecialCharTok{{-}}\NormalTok{ Root Directory Structures}\SpecialCharTok{:}
\NormalTok{  ├── stu}\SpecialCharTok{/}                         \CommentTok{\# Contains the official class roster (Excel/CSV)}
\NormalTok{  └── stat1}\SpecialCharTok{{-}}\NormalTok{ex}\SpecialCharTok{/}
\NormalTok{      ├── ps[no]}\SpecialCharTok{/}                  \CommentTok{\# Contains the question/problem set files}
\NormalTok{      └── ex[no]}\SpecialCharTok{/}                  \CommentTok{\# Contains student assignment submissions}
\end{Highlighting}
\end{Shaded}

\section{Execution Steps \& Workflow}\label{execution-steps-workflow}

\subsection{Step 1: Initialization \& Roster
Check}\label{step-1-initialization-roster-check}

\begin{enumerate}
\def\labelenumi{\arabic{enumi}.}
\tightlist
\item
  Access the official class roster in the \texttt{stu/} folder to
  retrieve the full list of student IDs and names.
\item
  Cross-reference the submissions in \texttt{stat1-ex/ex{[}no{]}/}.

  \begin{itemize}
  \tightlist
  \item
    If a student on the roster has not submitted a file, assign
    \textbf{Score: 0} and \textbf{Note: ``缺交'' (Missing)}.
  \end{itemize}
\end{enumerate}

\subsection{Step 2: Compliance \& Integrity Filtering
(Pre-grading)}\label{step-2-compliance-integrity-filtering-pre-grading}

For each submitted file, apply the following conditional checks
sequentially. If a file triggers a rule, stop further evaluation for
that file, assign \textbf{Score: 0}, and apply the designated note:

\begin{enumerate}
\def\labelenumi{\arabic{enumi}.}
\tightlist
\item
  \textbf{Deadline Check:}

  \begin{itemize}
  \tightlist
  \item
    \textbf{IF} the file submission timestamp is later than
    \texttt{2026-06-12\ 23:59\ (UTC+8)}
  \item
    \textbf{THEN} Score = 0, Note = ``Late Submission''
  \end{itemize}
\item
  \textbf{File Format \& Naming Check:}

  \begin{itemize}
  \tightlist
  \item
    \textbf{IF} the file is not a PDF OR does not follow the exact
    naming convention
    \texttt{ex{[}no{]}-{[}Student\_Name{]}-{[}Student\_ID{]}.pdf}
  \item
    \textbf{THEN} Score = 0, Note = ``Incorrect File Name/Format''
  \end{itemize}
\item
  \textbf{Internal YAML \& Code Metadata Check:}

  \begin{itemize}
  \item
    Open and inspect the source components of the submission.
  \item
    \textbf{IF} the document lacks the following YAML header elements or
    R code chunk:

\begin{Shaded}
\begin{Highlighting}[]
\FunctionTok{title}\KeywordTok{:}\AttributeTok{ }\StringTok{"Ex. [no]"}
\FunctionTok{author}\KeywordTok{:}\AttributeTok{ }\StringTok{"[Student\_Name] ([Student\_ID])"}
\FunctionTok{date}\KeywordTok{:}\AttributeTok{ }\StringTok{"YYYY{-}MM{-}DD"}
\end{Highlighting}
\end{Shaded}

    And the exact R execution chunk:

\begin{Shaded}
\begin{Highlighting}[]

\StringTok{\textasciigrave{}\textasciigrave{}\textasciigrave{}}\AttributeTok{ r}
\AttributeTok{print(uuid::UUIDgenerate())}
\end{Highlighting}
\end{Shaded}

\begin{verbatim}
## [1] "005b8221-4b15-4444-916b-5a91b913487c"
\end{verbatim}
  \item
    \textbf{THEN} Score = 0, Note = ``Incorrect Internal Format''
  \end{itemize}
\item
  \textbf{Plagiarism \& Academic Integrity Check:}

  \begin{itemize}
  \tightlist
  \item
    \textbf{IF} the generated string from \texttt{uuid::UUIDgenerate()}
    in a student's file is identical to another student's UUID
  \item
    \textbf{THEN} both/all involved students receive Score = 0, Note =
    ``Plagiarism''
  \end{itemize}
\end{enumerate}

\begin{center}\rule{0.5\linewidth}{0.5pt}\end{center}

\subsection{Step 3: Academic Content
Grading}\label{step-3-academic-content-grading}

For submissions that pass all filters in Step 2, grade the content
against the question set in \texttt{stat1-ex/ps{[}no{]}/}: -
\textbf{Score Calculation:} Total score is 100 points. Let \(N\) be the
total number of questions in the problem set. Each question is worth
\(\frac{100}{N}\) points (e.g., 50 points each for 2 questions; 25
points each for 4 questions). - \textbf{Partial Credit:} Award partial
credit based on accuracy and methodology. - \textbf{Deduction Note:} For
any incorrect answer, append the specific question number to the note
(Format: \texttt{{[}Question\_Number{]}is\ incorrect}).

\begin{center}\rule{0.5\linewidth}{0.5pt}\end{center}

\subsection{Step 4: Output Generation \& Statistical
Reporting}\label{step-4-output-generation-statistical-reporting}

\begin{enumerate}
\def\labelenumi{\arabic{enumi}.}
\tightlist
\item
  \textbf{Scoreboard Export (\texttt{ex-score.xlsx}):}

  \begin{itemize}
  \tightlist
  \item
    Create a new directory: \texttt{stu/ex{[}no{]}/}.
  \item
    Generate an Excel spreadsheet named \texttt{ex-score.xlsx} inside
    that directory.
  \item
    The file must contain exactly 4 columns: \texttt{StudentID},
    \texttt{Name}, \texttt{Score}, and \texttt{Note}.
  \end{itemize}
\item
  \textbf{Automated RMarkdown Analytics (\texttt{stat.Rmd} \&
  \texttt{stat.pdf}):}

  \begin{itemize}
  \tightlist
  \item
    In the \texttt{stu/ex{[}no{]}/} directory, dynamically generate an
    RMarkdown file named \texttt{stat.Rmd}.
  \item
    The \texttt{stat.Rmd} must read the data from \texttt{ex-score.xlsx}
    and perform the following:

    \begin{itemize}
    \tightlist
    \item
      \textbf{Descriptive Statistics:} Calculate minimum, \(Q_1\),
      median, \(Q_3\), maximum, and mean of the scores.
    \item
      \textbf{Data Visualization:} Render a clean boxplot of the score
      distribution.
    \item
      \textbf{Submissions Summary:} Calculate counts for Total
      Submissions, Missing, Late, and Plagiarism.
    \end{itemize}
  \item
    Knit/compile \texttt{stat.Rmd} into a production-ready
    \texttt{stat.pdf} in the same folder.
  \end{itemize}
\end{enumerate}

\end{document}
